% ===================================================================
%  TABEL Nr. 14 — De straat. — De winkeliers.
%  Traduction 2026 d'apres texto/io, controlee sur texto/fr et
%  texto/es. Consigne : voir texto/nl/10-tabel-01.tex.
% ===================================================================

%%K t14-apar-1 apar
\VUsaut{4.0mm}
\VUtitre{15.0pt}{60}{TABEL Nr. 14}
\VUsaut{3.0mm}

%%K t14-tit-1 sub
\VUpk{11.4pt}{40}{De straat. --- De winkeliers.}
\VUsaut{3.0mm}

%%K t14-01-1 p
1. --- Mijn vriend Jan, die op het land woont, kwam drie dagen bij mijn familie
doorbrengen. \VUgras{Tot} dan toe had hij geen gelegenheid gehad een grote stad
te zien; daarom gebruiken wij al onze dagen om te wandelen, en ik leg hem uit wat
hij niet kent.

%%K t14-02-1 p
2. --- Eerst gingen wij het \VUgras{observatorium} \textsuperscript{(1)}
bezoeken, dat hem zeer boeide. Toen wij terugkwamen door de \VUgras{straat}
\textsuperscript{(3)} waar het \VUgras{Turkse bad} \textsuperscript{(2)} is,
ontmoetten wij een \VUgras{begrafenis} \textsuperscript{(4)}. Voor de
\VUgras{lijkstoet} liep de \VUgras{geestelijkheid}, die de \VUgras{lijkwagen}
voorafging, waarin de \VUgras{kist} geplaatst was. Veel vrienden vergezelden de
\VUgras{rouwende} familie.

%%K t14-02-2 p
Nadat de stoet voorbij was, staken wij de straat over voor het grote
\VUgras{bierhuis} \textsuperscript{(5)} waarin veel \VUgras{gasten}
\VUgras{bier} dronken op het terras, beschut door de glazen \VUgras{markies}
\textsuperscript{(6)}.

%%K t14-03-1 p
3. --- Jan was verbaasd over het \VUgras{wapenschild} dat tussen de winkel van de
\VUgras{likeurstoker} \textsuperscript{(7)} en die van de \VUgras{muziekhandelaar}
\textsuperscript{(10)} hing. Hij vroeg mij naar de betekenis ervan; ik
antwoordde dat het het Engelse \VUgras{consulaat} \textsuperscript{(8)} aanduidt,
en wees hem op de \VUgras{consul} \textsuperscript{(9)} die op het balkon staat.

%%K t14-tit-2 sub
\VUpk{10.2pt}{40}{Verkenning van de winkels.}
\VUsaut{3.0mm}

%%K t14-04-1 p
4. --- Natuurlijk hielden wij stil voor de uitstalling van de
\VUgras{muziekinstrumenten}, \VUgras{harmoniums}, \VUgras{orgels} en
\VUgras{piano's}; wij bekeken de geïllustreerde omslagen van de
\VUgras{partituren} en van de \VUgras{muziekstukken} voor piano, en bewonderden
de \VUgras{lier} van verguld hout, die als uithangbord dient.

%%K t14-05-1 p
5. --- De stofwolken die de \VUgras{straatvegers} \textsuperscript{(11)} en de
\VUgras{veegwagen} \textsuperscript{(12)} maakten dwongen ons snel voorbij de
mooie \VUgras{meubelstellen} van de \VUgras{meubelhandelaar}
\textsuperscript{(13)} te lopen en voorbij de uitstalling van \VUgras{kachels}
en \VUgras{lampen} van de \VUgras{blikslager} \textsuperscript{(14)}.

%%K t14-06-1 p
6. --- Overigens waren wij op die plaats gedwongen het trottoir te verlaten, want
het was versperd door pakken, die men uit een \VUgras{verhuiswagen}
\textsuperscript{(16)} nam om ze in een \VUgras{pakhuis} \textsuperscript{(15)}
te zetten dat men zojuist verhuurd had.

%%K t14-06-2 p
Dat belette ons de mooie rijtuigen van de \VUgras{rijtuigmaker}
\textsuperscript{(17)} te zien, maar belette ons niet stil te houden om naar de
\VUgras{pakker} \textsuperscript{(19)} te kijken die een kist maakte voor zijn
buurman, de \VUgras{handelaar in fietsen en auto's} \textsuperscript{(18)}.
Daarna, omdat het al wat laat was, keerden wij naar huis terug.

%%K t14-tit-3 sub
\VUpk{10.2pt}{40}{Wandeling door de stad.}
\VUsaut{3.0mm}

%%K t14-07-1 p
7. --- De volgende dag stelde mijn moeder aan Jan voor dat men hem zou
fotograferen, daarna gaf zij mij de opdracht hem naar de \VUgras{fotograaf}
\textsuperscript{(20)} te vergezellen. Na het middagmaal gingen wij het
\VUgras{atelier} van de kunstenaar binnen, waar wij vrij lang moesten wachten.
Om ons bezig te houden bekeken wij de \VUgras{portretten}
\textsuperscript{(22)} die aan de muur hangen.

%%K t14-07-2 p
Toen het onze beurt was, duurde het werk niet lang, en zodra mijn vriend klaar
was met poseren voor het \VUgras{fototoestel} \textsuperscript{(21)}, daalden wij
af.

%%K t14-08-1 p
8. --- Op de eerste verdieping zagen wij door de deur van de \VUgras{kapper}
\textsuperscript{(23)}, die openstaat, zijn knecht die het haar van een klant
knipte.

%%K t14-08-2 p
Toen wij op straat gekomen waren, kwamen wij langs de \VUgras{tabakswinkel}
\textsuperscript{(24)}. Jan werd zo vermaakt door de rode \VUgras{lantaarn}
\textsuperscript{(25)} en de \VUgras{dikke pijp} die als \VUgras{uithangbord}
\textsuperscript{(26)} dient, dat hij tegen twee oude heren botste die pratend
\VUgras{snuiftabak namen} \textsuperscript{(27)}.

%%K t14-tit-4 sub
\VUpk{10.2pt}{40}{De banketbakkerij.}
\VUsaut{3.0mm}

%%K t14-09-1 p
9. --- De \VUgras{bankbiljetten} en \VUgras{muntstukken} uit alle landen die in
de \VUgras{etalage} van de \VUgras{wisselaar} \textsuperscript{(29)} uitgestald
waren trokken enkele ogenblikken onze aandacht, maar omdat het tijd was voor het
vieruurtje, gingen wij bij de \VUgras{banketbakker} \textsuperscript{(31)} naar
binnen zonder langer stil te houden.

%%K t14-10-1 p
10. --- Toen wij alle \VUgras{lekkernijen} zagen die de winkel bevatte,
\VUgras{taarten, ontbijtkoek, koekjes} en allerlei \VUgras{snoepgoed}, konden wij
niet gemakkelijk kiezen. Ten slotte koos Jan de \VUgras{roomtaartjes}, en ik nam
alleen een klein \VUgras{kersentaartje}, want van het zoet \VUgras{krijg ik
kiespijn}, en ik verlang er helemaal niet naar te dikwijls de \VUgras{tandarts}
\textsuperscript{(30)} te gaan bezoeken.

%%K t14-tit-5 sub
\VUpk{10.2pt}{40}{Machineschrijven.}
\VUsaut{3.0mm}

%%K t14-11-1 p
11. --- Om aan mijn vriend een \VUgras{schrijfmachine} te tonen waarover hij had
horen spreken, maar die hij niet kende, gingen wij het \VUgras{kantoor}
\textsuperscript{(32)} van mijn \VUgras{neef} \textsuperscript{(34)} binnen. Wij
troffen hem aan terwijl hij zijn brieven dicteerde aan zijn
\VUgras{secretaris} \textsuperscript{(35)} die \VUgras{stenograaf} is.
Nauwelijks was een brief af, of een \VUgras{typist} \textsuperscript{(33)}
schreef hem met zijn machine over. Omdat wij de bezigheden van mijn bloedverwant
niet wilden onderbreken, bleven wij maar enkele ogenblikken bij hem.

%%K t14-12-1 p
12. --- Onder het kantoor van mijn neef woont een grote
\VUgras{horlogemaker-juwelier} \textsuperscript{(37)} die bovendien goudsmid is.
Zijn prachtige winkel bevat veel \VUgras{klokken, pendules, zakhorloges},
\VUgras{kettinkjes} en kostbare \VUgras{juwelen}, bijvoorbeeld
\VUgras{parelsnoeren}, gouden \VUgras{armbanden}, \VUgras{oorbellen} en
\VUgras{ringen} versierd met \VUgras{edelstenen} en \VUgras{diamanten}. Een van
de etalages is bijzonder gereserveerd voor het \VUgras{goudwerk} en bevat een
belangrijke verzameling zilveren \VUgras{eetgerei}.

%%K t14-13-1 p
13. --- Toen wij de straathoek omsloegen, merkte ik dat men het \VUgras{bordje}
\textsuperscript{(36)} verwisseld had dat bij het \VUgras{huisnummer} zit. Ik
stond op het punt mijn opmerking hardop te zeggen, toen de vrolijke klanken van
\VUgras{de} militaire muziek gehoord werden. Wij begonnen het regiment tegemoet
te rennen, zonder te bedenken dat \VUgras{het} voorbij de winkels van de
\VUgras{hoedenmaker} \textsuperscript{(39)} en van de \VUgras{handschoenmaker}
\textsuperscript{(40)} zou komen, en dat wij even goed geplaatst zouden zijn
geweest om zijn defilé te zien als wij voor de etalage van de \VUgras{opticien}
\textsuperscript{(38)} gebleven waren, die geheel voorzien is van
\VUgras{knijpbrillen, lorgnetten} en \VUgras{verrekijkers}.

%%K t14-tit-6 sub
\VUpk{10.2pt}{40}{Marsdefilé.}
\VUsaut{3.0mm}

%%K t14-14-1 p
14. --- Voor het \VUgras{regiment} \textsuperscript{(41)} liep de
\VUgras{tamboer-majoor} \textsuperscript{(42)} gevolgd door de
\VUgras{tamboers} \textsuperscript{(43)}, de \VUgras{hoornblazers}
\textsuperscript{(44)} en het hele \VUgras{muziekkorps}. De \VUgras{generaal}
\textsuperscript{(45)} die met zijn adjudant op het plein stond, was blijven
staan bij de \VUgras{waterstraal} \textsuperscript{(49)} van de
\VUgras{fontein} \textsuperscript{(48)}, om het \VUgras{defilé} bij te wonen.

%%K t14-14-2 p
\VUgras{De stafofficieren} \textsuperscript{(46)}, de \VUgras{kolonel} en de
\VUgras{luitenant-kolonel} groetten hem met de degen. \VUgras{De majoors}
stonden voor hun \VUgras{bataljons} \textsuperscript{(47)} en elke
\VUgras{kapitein} voerde zijn \VUgras{compagnie} aan, terwijl de
\VUgras{luitenants}, \VUgras{onderluitenants} en \VUgras{onderofficieren} hun
gewone plaats bij hun mannen innamen.

%%K t14-15-1 p
15. --- Veel voorbijgangers hadden zich op \VUgras{het kruispunt}
\textsuperscript{(50)} verzameld; de winkeliers, de \VUgras{paraplumaker}
\textsuperscript{(51)}, de \VUgras{verver} \textsuperscript{(53)}, de
\VUgras{wapensmid} \textsuperscript{(54)} kwamen naar buiten om de
\VUgras{soldaten} te zien. Alle voertuigen, \VUgras{huurrijtuigen}
\textsuperscript{(52)}, \VUgras{herenrijtuigen} \textsuperscript{(57)} en ook de
\VUgras{sproeiwagen} \textsuperscript{(56)}, gingen langs het trottoir staan.

%%K t14-tit-7 sub
\VUpk{10.2pt}{40}{Straattaferelen.}
\VUsaut{3.0mm}

%%K t14-15-2 p
Een \VUgras{handelsreiziger} \textsuperscript{(55)} die zijn
\VUgras{monsterkoffers} in de hand droeg, stak snel de straat over, en de mensen
die op de \VUgras{imperiaal} \textsuperscript{(59)} van de \VUgras{omnibus}
\textsuperscript{(58)} zaten keerden zich om om van het schouwspel te genieten.
Een arme \VUgras{straatzanger} \textsuperscript{(60)} met zijn \VUgras{draaiorgel}
\textsuperscript{(61)} moest zijn \VUgras{lied} onderbreken, want het lawaai van
de tamboers overstemde zijn stem.

%%K t14-16-1 p
16. --- Toen het regiment voor de \VUgras{stoffenwinkel} \textsuperscript{(64)}
kwam verlieten de \VUgras{naaisters} \textsuperscript{(63)} en de
\VUgras{kleermakers} \textsuperscript{(62)} hun \VUgras{naaimachines} en hun
werk om door het raam te kijken. \VUgras{De winkelier} \textsuperscript{(65)},
die zojuist nieuwe \VUgras{kostuums} \textsuperscript{(66)} in zijn
\VUgras{uitstalling} gezet had, moest de \VUgras{lakens} \textsuperscript{(67)}
en het \VUgras{linnen} naar binnen laten brengen die op het trottoir stonden,
bij de \VUgras{rioolopening} \textsuperscript{(68)}, uit vrees dat de menigte ze
zou omgooien, zelfs een oude \VUgras{marskramer} \textsuperscript{(69)}, van wie
een portierster kousen afdong, was gedwongen een gang binnen te gaan om zijn
verkoop af te maken.

%%K t14-16-2 p
Wij vergezelden het regiment tot aan de kazerne, \VUgras{en}, zeer tevreden over
onze dag, keerden wij op het uur van het diner terug.

%%K t14-tit-8 sub
\VUpk{10.2pt}{40}{Verschillende bedrijven en beroepen.}
\VUsaut{3.0mm}

%%K t14-17-1 p
17. --- De volgende dag stelde mijn vader ons voor de drukkerij van de
plaatselijke krant te bezoeken. Wij namen dat geestdriftig aan.

%%K t14-17-2 p
De \VUgras{drukker} is ook \VUgras{boekhandelaar} \textsuperscript{(70)}
\VUgras{(= boekverkoper)}, \VUgras{uitgever}, \VUgras{papierhandelaar} en
\VUgras{boekbinder}. Hij heeft op de eerste verdieping de \VUgras{redactie}
\textsuperscript{(71)} van de krant ondergebracht, en ook het
\VUgras{graveeratelier}, waarin de \VUgras{lithografen} \textsuperscript{(72)} en
de \VUgras{kopergraveurs} \textsuperscript{(73)} werken, ten slotte de
\VUgras{zetterij}, waarin de \VUgras{drukkersgezellen} \textsuperscript{(74)}
zijn. De eigenlijke \VUgras{drukkerij} \textsuperscript{(75)}, de
\VUgras{drukpersen} en de verschillende machines staan op de begane grond.

%%K t14-tit-9 sub
\VUpk{10.2pt}{40}{Jan stelt zeer veel vragen.}
\VUsaut{3.0mm}

%%K t14-18-1 p
18. --- Toen wij de drukkerij uitkwamen, hield Jan zeer verbaasd stil voor een
klein glazen kastje dat op een zuil bevestigd was, tegenover de
\VUgras{garenwinkel} \textsuperscript{(77)} en vroeg waarvoor het dient. Mijn
vader legde hem uit, dat het een \VUgras{brandmelder} \textsuperscript{(76)} is.
Overigens was alles, in de stad, voor mijn vriend een wonder, van de eenvoudige
\VUgras{stenen fontein} \textsuperscript{(78)} en de lichte
\VUgras{bestellingsdriewieler} \textsuperscript{(80)} die door een
\VUgras{piccolo} \textsuperscript{(79)} in livrei gereden wordt, tot de
\VUgras{postwagen} \textsuperscript{(81)}.

%%K t14-19-1 p
19. --- Terwijl mijn vader ons enkele ogenblikken verliet om met de
\VUgras{belastingontvanger} \textsuperscript{(83)} te praten, die was blijven
staan om met een \VUgras{advocaat} \textsuperscript{(82)} en een
\VUgras{notaris} \textsuperscript{(84)} te spreken, vermaakten wij ons door het
verkeer op straat met de ogen te volgen. De meest uiteenlopende mensen kwamen
langs ons : een \VUgras{banketbakkersknecht} \textsuperscript{(85)}, die in een
mand een prachtige \VUgras{pastei} droeg, kwam achter een
\VUgras{kleerkoper} \textsuperscript{(86)}; een \VUgras{lantaarnopsteker}
\textsuperscript{(87)} praatte met een \VUgras{gasfitter}
\textsuperscript{(88)}, een \VUgras{non} \textsuperscript{(89)} die een weesmeisje
bij de hand hield keerde naar haar klooster terug.

%%K t14-tit-10 sub
\VUpk{10.2pt}{40}{Op onze terugweg naar huis.}
\VUsaut{3.0mm}

%%K t14-20-1 p
20. --- Op het ogenblik dat mijn vader zich weer bij ons voegde, lachte Jan
luidkeels, toen hij het vuile gezicht en de vreemde haveloze kleding zag van twee
\VUgras{schoorsteenvegers} \textsuperscript{(91)} die helemaal met roet bedekt
waren.

%%K t14-21-1 p
21. --- Ik wilde bij de \VUgras{bloemenverkoopster} \textsuperscript{(92)} een
plant voor mijn moeder kopen, maar mijn vader wees mij erop, dat die zeer zwaar
zou zijn om mee te dragen en dat het beter zou zijn \VUgras{sinaasappels} te
kopen. Wij naderden de \VUgras{verkoopster} \textsuperscript{(94)} en toen de
\VUgras{gouvernante} \textsuperscript{(93)}, die voor haar pupil aan het kiezen
was, haar aankoop klaar had, lieten wij ons bedienen.

%%K t14-22-1 p
22. --- Toen wij naar huis terugkeerden, ontmoetten wij verscheidene kennissen :
een oude vriendin van mijn familie, die op de arm van haar \VUgras{gezelschapsdame}
\textsuperscript{(95)} steunde, de \VUgras{ingenieur} \textsuperscript{(96)} van
bruggen en wegen, en een van mijn nichten met haar \VUgras{kindermeisje}
\textsuperscript{(98)}.

%%K t14-22-2 p
Daarna ontmoetten wij de \VUgras{modiste} \textsuperscript{(97)} van mijn moeder
en twee \VUgras{monniken} \textsuperscript{(99)} die vreemd aan de stad waren,
die naar ons toe kwamen om mijn vader naar de weg naar het station te vragen.
\VUsaut{6.0mm}
\VUfilet{22mm}
